\documentclass[../../main.tex]{subfiles}
\begin{document}

{\bf [解]}

\paragraph{【新課程】}

実数$\theta\, (0\le\theta <2\pi)$について簡単のため$\cos\theta = c$，$s = \sin\theta$とおく．すると円
\begin{align}
  \left(x-5\right)^2 + \left(y-5\right)^2 = r^2 \label{eq:1}
\end{align}
の上の点$P$は$P(5+rc, 5+rs)$と書ける．

\begin{figure}[htb]
  \centering
  \begin{tikzpicture}[scale=0.35, >=stealth, every node/.style={font=\small}]

    % --- 1. 定数・中心座標の定義 ---
    % 円 C1 の中心 O1(5, 5), 半径 r = 3 (視認性重視)
    \def\rVal{3.0}
    \def\tVal{40} % パラメータ theta = 40°

    \coordinate (O1) at (5, 5);
    \coordinate (A)  at (9, 0);

    % P の座標: (5 + r*cos t, 5 + r*sin t)
    \coordinate (P) at ($ (O1) + ({\tVal}:{\rVal}) $);

    % Q: A(9,0) に関する P の対称点 (中心 O2(13, -5))
    \coordinate (O2) at (13, -5);
    \coordinate (Q) at ($ (O2) + ({180+\tVal}:{\rVal}) $);

    % R: P を原点まわりに +90° 回転 (中心 O3(-5, 5))
    \coordinate (O3) at (-5, 5);
    \coordinate (R) at ($ (O3) + ({90+\tVal}:{\rVal}) $);

    % --- 2. 座標軸 ---
    \draw[->, thick] (-11, 0) -- (19, 0) node[right] {$x$};
    \draw[->, thick] (0, -10) -- (0, 11) node[above] {$y$};
    \node[below left=2pt] at (0, 0) {$O$};

    % --- 3. 3つの円 (C1: P の軌跡, C2: Q の軌跡, C3: R の軌跡) ---
    % 元の円 C1
    \draw[very thick, blue!80!black] (O1) circle (\rVal);
    \node[blue!80!black, above right] at ($ (O1) + (45:\rVal) $) {$C_1$};

    % Q の軌跡 C2 (中心 (13, -5))
    \draw[thick, dashed, purple!80!black] (O2) circle (\rVal);
    \node[purple!80!black, below right] at ($ (O2) + (-45:\rVal) $) {$C_2$ ($Q$ の軌跡)};

    % R の軌跡 C3 (中心 (-5, 5))
    \draw[thick, dashed, teal!80!black] (O3) circle (\rVal);
    \node[teal!80!black, above left] at ($ (O3) + (135:\rVal) $) {$C_3$ ($R$ の軌跡)};

    % --- 4. 幾何変換の補助線 ---
    % (1) P と Q の点対称変換 (点 A に関する対称)
    \draw[thick, dashed, gray!70] (P) -- (Q);
    \fill[red!80!black] (A) circle (3pt) node[below=3pt] {$A(9,\, 0)$};

    % (2) P から R への +90° 回転
    \draw[thick, dashed, gray!60] (0, 0) -- (P);
    \draw[thick, dashed, gray!60] (0, 0) -- (R);
    \draw[->, thick, teal!80!black] (2.2, 2.2) arc[start angle=45, end angle=135, radius=3.1];
    \node[teal!80!black, above=2pt, font=\footnotesize] at (0, 3.2) {$+90^\circ$ 回転};

    % (3) 線分 QR (|QR| の最小・最大を考える対象)
    \draw[very thick, red!85!black] (Q) -- (R)
    node[midway, fill=white, inner sep=1.5pt, font=\normalsize\bfseries] {$|QR|$};

    % --- 5. 各点と中心のプロットとラベル ---
    % 中心
    \fill (O1) circle (2.5pt) node[above right=1pt, font=\footnotesize] {$(5,5)$};
    \fill (O2) circle (2.5pt) node[below left=1pt, font=\footnotesize] {$(13,-5)$};
    \fill (O3) circle (2.5pt) node[above left=1pt, font=\footnotesize] {$(-5,5)$};

    % 点 P, Q, R
    \fill[blue!80!black] (P) circle (3.5pt) node[above right] {$P$};
    \fill[purple!80!black] (Q) circle (3.5pt) node[below left] {$Q$};
    \fill[teal!80!black] (R) circle (3.5pt) node[above left] {$R$};

    % 各中心間を結ぶ基準線（中心間距離 \sqrt{18^2 + 10^2} = 2\sqrt{53}）
    \draw[dotted, gray!80!black, thick] (O3) -- (O2)
    node[pos=0.25, below left, font=\footnotesize] {$|O_3 O_2| = 2\sqrt{53}$};

  \end{tikzpicture}
  \caption{円 $C_1$ 上の点 $P$、点 $A$ に関する対称点 $Q$、原点中心 $+90^\circ$ 回転点 $R$ の幾何配置}
  \label{fig:p_q_r_geometry}
\end{figure}

従って $A(9,0)$ に関する$P$の対称点は
\begin{align}
  Q(13-rc, -5-rs) \label{eq:2}
\end{align}
である．

$R$を正の方向に$\dfrac{\pi}{2}$回転させた$R$は
\begin{align}
  R\left(-5-rs,5+rc\right) \label{eq:3}
\end{align}
である．

従って\cref{eq:2,eq:3}から$|QR|^2$は
\begin{align}
  |QR|^2
  &= \{(13-rc) - (-5-rs)\}^2 + \{(-5-rs) - (5+rc)\}^2 \\
  &= (18-rc+rs)^2 + (-10-rs-rc)^2 \\
  &= 18^2 + 10^2 + 2r^2 + 56rs - 16rc \\
  &= 424 + 2r^2 + 8\sqrt{53}r\sin(\theta-\alpha) \label{eq:4}
\end{align}
とかける．ここで $\alpha$ は
\begin{align}
  \tan\alpha = \dfrac{2}{7} \label{eq:5}
\end{align}
を満たす数である．$0 \leqq \theta < 2\pi$から，
\begin{align}
  -\alpha \leqq \theta - \alpha < 2\pi - \alpha
\end{align}
であるから，$-1 \leqq \sin(\theta-\alpha) \leqq 1$である．
ゆえに$r>0$から，$|QR|^2$の最小最大値$f(r),g(r)$は
\begin{align}
  \begin{cases}
    f(r) = \sqrt{424-8\sqrt{53}r+2r^2} = \sqrt{2}|r-2\sqrt{53}| \\
    g(r) = \sqrt{424+8\sqrt{53}r+2r^2} = \sqrt{2}(r+2\sqrt{53})
  \end{cases}
\end{align}
とかける．
また$f(r)=0$のときは$r=2\sqrt{53}$である．$\cdots\text{（答）}$

\paragraph{【旧課程】}

\begin{figure}[htb]
  \centering
  \begin{tikzpicture}[scale=0.75, >=stealth, every node/.style={font=\small}]

    % --- 座標軸 ---
    \draw[->, thick, gray!60] (-4.5, 0) -- (10.0, 0) node[right] {$X$};
    \draw[->, thick, gray!60] (0, -5.5) -- (0, 7.5) node[above] {$Y$};
    \node[below left=2pt] at (0, 0) {$O$};

    % --- 各円の中心と半径 ---
    \coordinate (CenterA) at (0, 0);
    \coordinate (CenterB) at (4, 3);
    \coordinate (CenterC) at (5, -3);

    \def\rA{3.0}
    \def\rB{2.0}
    \def\rC{1.0}

    % 円の描画
    \draw[very thick, blue!75!black] (CenterA) circle (\rA);
    \node[blue!75!black, left] at (-3, 0.5) {円 $A$};

    \draw[very thick, orange!85!black] (CenterB) circle (\rB);
    \node[orange!85!black, right] at (4, 5.5) {円 $B$};

    \draw[very thick, teal!85!black] (CenterC) circle (\rC);
    \node[teal!85!black, right] at (6, -3) {円 $C$};

    % 各円の中心
    \fill (CenterA) circle (1.8pt);
    \fill (CenterB) circle (1.8pt) node[below=2pt] {$(4,3)$};
    \fill (CenterC) circle (1.8pt) node[below=2pt] {$(5,-3)$};

    % --- 点 P の配置 (軌跡 Gamma 上の一例として theta = 50° を採用) ---
    % Gamma: (X-3)^2 + Y^2 = 27 (半径 3*sqrt(3) ≈ 5.196)
    \coordinate (M) at (3, 0);
    \def\rGamma{5.19615}
    \coordinate (P) at ($ (M) + (50:\rGamma) $); % 概ね (6.34, 3.98)

    \fill[purple!80!black] (P) circle (2.2pt) node[above right] {$P(X,Y)$};

    % --- 接線と接点 (幾何計算) ---
    % P から円 A への接点 TA
    \pgfmathsetmacro{\distPA}{sqrt((6.34)^2 + (3.98)^2)} % ≈ 7.487
    \pgfmathsetmacro{\angPA}{atan2(3.98, 6.34)}
    \pgfmathsetmacro{\dAngA}{acos(\rA / \distPA)}
    \coordinate (TA) at ($ (CenterA) + ({\angPA + \dAngA}:\rA) $);

    % P から円 B への接点 TB
    \pgfmathsetmacro{\distPB}{sqrt((6.34 - 4)^2 + (3.98 - 3)^2)} % ≈ 2.541
    \pgfmathsetmacro{\angPB}{atan2(3.98 - 3, 6.34 - 4)}
    \pgfmathsetmacro{\dAngB}{acos(\rB / \distPB)}
    \coordinate (TB) at ($ (CenterB) + ({\angPB - \dAngB}:\rB) $);

    % P から円 C への接点 TC
    \pgfmathsetmacro{\distPC}{sqrt((6.34 - 5)^2 + (3.98 - (-3))^2)} % ≈ 7.107
    \pgfmathsetmacro{\angPC}{atan2(3.98 - (-3), 6.34 - 5)}
    \pgfmathsetmacro{\dAngC}{acos(\rC / \distPC)}
    \coordinate (TC) at ($ (CenterC) + ({\angPC + \dAngC}:\rC) $);

    % 接線の描画
    \draw[thick, blue!80!black] (P) -- (TA) node[midway, above left, font=\footnotesize] {$\alpha(P)$};
    \draw[thick, orange!90!black] (P) -- (TB) node[midway, right=1pt, font=\footnotesize] {$\beta(P)$};
    \draw[thick, teal!80!black] (P) -- (TC) node[midway, right=1pt, font=\footnotesize] {$\gamma(P)$};

    % 接点のプロットと半径（直角記号用点線）
    \fill (TA) circle (1.4pt) node[above left, font=\scriptsize] {$T_A$};
    \fill (TB) circle (1.4pt) node[below=2pt, font=\scriptsize] {$T_B$};
    \fill (TC) circle (1.4pt) node[below left, font=\scriptsize] {$T_C$};

    \draw[dashed, gray!70] (CenterA) -- (TA);
    \draw[dashed, gray!70] (CenterB) -- (TB);
    \draw[dashed, gray!70] (CenterC) -- (TC);

    % 条件式の注記
    \node[draw=purple!80!black, rounded corners=3pt, fill=white, inner sep=4pt, font=\footnotesize]
    at (1.0, -4.5) {$\alpha(P)^2 + \beta(P)^2 + \gamma(P)^2 = 99$};

  \end{tikzpicture}
  \caption{点 $P$ から 3 つの円 $A, B, C$ への接線とその距離}
  \label{fig:three_circles_tangents_concept}
\end{figure}

点$P$の座標を$(X, Y)$とする．
中心$K$，半径$R$の円に対して点$P$からその円に接線が引けるための必要十分条件は，点$P$が円の外部または周上にあること，すなわち$|PK|^2 \geqq R^2$である．
このとき，接点までの距離の$2$乗は三平方の定理より$|PK|^2 - R^2$に等しい．
したがって，各円について条件は以下の通りとなる．

\subparagraph{円$A$}
中心$(0, 0)$，半径$3$であるから，接線が引ける条件は
\begin{align}
  X^2 + Y^2 \geqq 9 \label{eq:5}
\end{align}
であり，距離は
\begin{align}
  \alpha(P)^2 = X^2 + Y^2 - 9 \label{eq:6}
\end{align}
である．

\subparagraph{円$B$}
中心$(4, 3)$，半径$2$であるから，接線が引ける条件は
\begin{align}
  (X-4)^2 + (Y-3)^2 \geqq 4 \label{eq:7}
\end{align}
であり，距離は
\begin{align}
  \beta(P)^2 = (X-4)^2 + (Y-3)^2 - 4 \label{eq:8}
\end{align}
である．

\subparagraph{円$C$}
中心$(5, -3)$，半径$1$であるから，接線が引ける条件は
\begin{align}
  (X-5)^2 + (Y+3)^2 \geqq 1 \label{eq:9}
\end{align}
であり，距離は
\begin{align}
  \gamma(P)^2 = (X-5)^2 + (Y+3)^2 - 1 \label{eq:10}
\end{align}
である．
\casesend

\cref{eq:6,eq:8,eq:10}を題意の条件式$\alpha(P)^2 + \beta(P)^2 + \gamma(P)^2 = 99$に代入して整理すると，
\begin{align}
  &(X^2 + Y^2 - 9) + \{(X-4)^2 + (Y-3)^2 - 4\} + \{(X-5)^2 + (Y+3)^2 - 1\} = 99 \nonumber\\
  &\Longleftrightarrow 3X^2 + 3Y^2 - 18X + 45 = 99 \nonumber\\
  &\Longleftrightarrow X^2 - 6X + Y^2 = 18 \nonumber\\
  &\Longleftrightarrow (X-3)^2 + Y^2 = 27 \label{eq:11}
\end{align}
を得る．この円を$\Gamma$とし，中心を$M(3, 0)$，半径を$r = 3\sqrt{3}$とおく．
従って求める点$P$の条件は$\Gamma$の円周上で\cref{eq:5,eq:7,eq:9}を満たす部分である．それぞれの条件について個別に確認する．

\subparagraph{円$A$: \cref{eq:5}}

円$\Gamma$上の点を$P(X,Y)=(r\cos\theta+3,r\sin\theta)$とおくと円$A$の中心$(0,0)$との距離の二乗は
\begin{align}
  X^2 + Y^2
  &= (r\cos\theta+3)^2+\left(r\sin\theta\right)^2 \\
  &= r^2 + 9 + 6r\cos\theta \\
\end{align}
であり，これが円$A$の半径の二乗$9$より大きい条件は
\begin{align}
  r^2 + 9 + 6r\cos\theta \ge 9 \iff r^2 + 6r\cos\theta \ge 0 \iff \cos\theta \ge -\dfrac{r}{6} = -\dfrac{\sqrt{3}}{2}
\end{align}
であり，従って円$A$の外側の条件は
\begin{align}
  0 \le \theta \dfrac{5\pi}{6}, \dfrac{7\pi}{6}\le \theta < 2\pi \label{eq:12}
\end{align}
である．

\subparagraph{円$B$: \cref{eq:7}}

円$\Gamma$上の点を$P(X,Y)=(r\cos\theta+3,r\sin\theta)$とおくと円$B$の中心$(4,3)$との距離の二乗は
\begin{align}
  (X-4)^2 + (Y-3)^2
  &= (r\cos\theta+3-4)^2+\left(r\sin\theta-3\right)^2 \\
  &= r^2 + 10 - 2r(\cos\theta + 4\sin\theta) \\
  &\ge r^2 + 10 - 2r\sqrt{5} \\
  &= 27 + 10 -6\sqrt{15} \\
  &= 37 - 6\sqrt{15} \ge 4
\end{align}
となり，\cref{eq:7}を満たす．ただし3行目で三角関数の合成を用いた．
よって円$\Gamma$上の点は全て円$B$の外側にある．

\subparagraph{円$C$: \cref{eq:9}}

円$\Gamma$上の点を$P(X,Y)=(r\cos\theta+3,r\sin\theta)$とおくと円$C$の中心$(5,-3)$との距離の二乗は
\begin{align}
  (X-5)^2 + (Y+3)^2
  &= (r\cos\theta+3-5)^2+\left(r\sin\theta+3\right)^2 \\
  &= r^2 + 13 + 2r(-2\cos\theta + 3\sin\theta) \\
  &\ge r^2 + 13 - 2r\sqrt{13} \\
  &= 27 + 13 -6\sqrt{39} \\
  &= 40 - 6\sqrt{39} \ge 1
\end{align}
となり，\cref{eq:9}を満たす．ただし3行目で三角関数の合成を用いた．
よって円$\Gamma$上の点は全て円$C$の外側にある．
\casesend

以上全ての場合分けから，求める条件は\cref{eq:12}である．
この領域を図示すると下図の太線部である（円の内部に含まれる破線部は除く）．

\begin{figure}[htb]
  \centering
  \begin{tikzpicture}[scale=0.85, >=stealth, every node/.style={font=\small}]

    % --- 1. 定数定義 ---
    % Gamma: 中心 M(3, 0), 半径 r = 3*sqrt(3) ≈ 5.196
    % 円 A: 中心 (0, 0), 半径 3
    % 円 A と Gamma の交点:
    % X = -1.5, Y = ± 3*sqrt(3)/2 ≈ ± 2.598 (偏角 θ = ±150°)
    \coordinate (O)  at (0, 0);
    \coordinate (M)  at (3, 0);
    \coordinate (B)  at (4, 3);
    \coordinate (C)  at (5, -3);

    \def\rGamma{5.19615} % 3*sqrt(3)
    \coordinate (P1) at (-1.5,  2.59808);
    \coordinate (P2) at (-1.5, -2.59808);

    % --- 2. 座標軸 ---
    \draw[->, thick] (-4.0, 0) -- (9.8, 0) node[right] {$X$};
    \draw[->, thick] (0, -5.5) -- (0, 6.2) node[above] {$Y$};
    \node[below left=2pt] at (O) {$O$};

    % --- 3. 円 A, B, C の描画 ---
    % 円 A (半径 3)
    \draw[thick, gray!70] (O) circle (3);
    \node[gray!80!black, left] at (-3, 0.4) {円 $A$};

    % 円 B (中心 (4,3), 半径 2)
    \draw[thick, gray!70] (B) circle (2);
    \fill[gray!70] (B) circle (1.2pt);
    \node[gray!80!black, above right] at (B) {円 $B(4,3)$};

    % 円 C (中心 (5,-3), 半径 1)
    \draw[thick, gray!70] (C) circle (1);
    \fill[gray!70] (C) circle (1.2pt);
    \node[gray!80!black, below right] at (C) {円 $C(5,-3)$};

    % --- 4. 境界直線 X = -3/2 ---
    \draw[thick, dashed, gray!60] (-1.5, -4.5) -- (-1.5, 4.5)
    node[above, font=\footnotesize, black] {$X = -\dfrac{3}{2}$};

    % --- 5. 点 P の軌跡 Gamma ---
    % 円 A の外部: 偏角 -150° から +150° (実線・太線)
    \draw[line width=1.6pt, blue!80!black]
    (M) ++(-150:\rGamma) arc[start angle=-150, end angle=150, radius=\rGamma];

    % 円 A の内部: 偏角 150° から 210° (-150°) (破線・除外部分)
    \draw[thick, dashed, red!70!black]
    (M) ++(150:\rGamma) arc[start angle=150, end angle=210, radius=\rGamma];

    % --- 6. 中心 M から交点 P1, P2 への動径補助線と角度 ---
    \draw[dashed, blue!50!black] (M) -- (P1);
    \draw[dashed, blue!50!black] (M) -- (P2);
    \draw[thick, blue!60!black] (M) +(0:0.9) arc[start angle=0, end angle=150, radius=0.9];
    \node[blue!70!black, font=\footnotesize] at ($(M) + (75:1.2)$) {$\dfrac{5\pi}{6}$};

    % --- 7. 主要点・ラベル ---
    % 中心 M
    \fill (M) circle (1.6pt) node[below=2pt] {$M(3,0)$};

    % 交点 P1, P2 (白抜き丸で除外境界を明示、または黒丸)
    \fill[purple!80!black] (P1) circle (1.8pt) node[above left] {$P_1\left(-\frac{3}{2},\, \frac{3\sqrt{3}}{2}\right)$};
    \fill[purple!80!black] (P2) circle (1.8pt) node[below left] {$P_2\left(-\frac{3}{2},\, -\frac{3\sqrt{3}}{2}\right)$};

    % 軌跡ラベル
    \node[blue!80!black, right, font=\bfseries] at (8.2, 0.4) {$\Gamma$};
    % \node[blue!80!black, fill=white, inner sep=2pt, font=\footnotesize]
    % at (6.2, 3.8) {軌跡（太線部 $L=5\sqrt{3}\pi$）};

  \end{tikzpicture}
  \caption{点 $P$ の軌跡 $\Gamma$（太線部）と除外される円弧（破線部）}
  \label{fig:locus_gamma_final}
\end{figure}

この円弧の中心角は
\begin{align}
  \Theta = 2\pi - \dfrac{2\pi}{6} = \dfrac{5\pi}{3}
\end{align}
であり，円の半径$r$だから円弧の長さは
\begin{align}
  L = r \dfrac{5\pi}{3} = 5\sqrt{3}\pi
\end{align}
である．

\end{document}
